\documentclass{article}

\usepackage{booktabs}
\usepackage{tabularx}

\title{Development Plan\\\progname}

\author{\authname}

\date{}

\input{../Comments.text}
\input{../Common.text}

\begin{document}

\maketitle

\begin{table}[hp]
\caption{Revision History} \label{TblRevisionHistory}
\begin{tabularx}{\textwidth}{llX}
\toprule
\textbf{Date} & \textbf{Developer(s)} & \textbf{Change}\\
\midrule
September 23, 2026 & Andreas Sideris & Initial draft of Workflow 
Plan\\
September 24, 2026 & Andreas Sideris & Initial draft of Project 
Decomposition and Scheduling\\
September 25, 2026 & Andreas Sideris & Use of AI and LLMs Completed\\
September 26, 2026 & Owen Johnson & Confidential Information and IP to Protect 
Completed\\
September 26, 2026 & Christopher Ludwig & Team Charter, Communication, 
Meetings, Roles Rough Draft\\
September 26, 2026 & Andreas Sideris & Coding Standard Added\\
September 27, 2026 & Owen Johnson & PoC Demo Plan Expanded, Expected Technology Slightly Edited\\
September 27, 2026 & Humna Saeed & Introduction, Team Charter\\ 
\bottomrule
\end{tabularx}
\end{table}

\newpage{}

The following development plan covers critical areas that will aid throughout the development of RealmForge, including core project management decisions, team organization, project scheduling, workflow, proof of concept, and expected technologies.

%\wss{Additional information on the development plan can be found in the
%\href{https://gitlab.cas.mcmaster.ca/courses/capstone/-/blob/main/Lectures/L02b_POCAndDevPlan/POCAndDevPlan.pdf?ref_type=heads}
%{lecture slides}.}

\section{Confidential Information?}

\textbf{Protected Information:} This project does not have any confidential information to protect.



\section{IP to Protect}

\textbf{Project Type:} Internal Capstone
\\

\noindent Based on the 
\href {https://gitlab.cas.mcmaster.ca/courses/capstone/-/blob/main/Lectures/L05_NDA_IP_Licenses/MILO_Templates/Capstone_Projects_IP_Guidelines.docx?ref_type=heads}{{Capstone Projects IP Guidelines}} document in the course repository, this project is an internal capstone. This means that the students own the IP by default, no agreement or acknowledgement is required, and MILO does not need to be involved. 



\section{Copyright License}

\wss{What copyright license is your team adopting.  Point to the license in your
repo.}

\section{Team Meeting Plan}

Meetings will be held every Monday from 2:30pm - 3:30pm. More meetings will be 
scheduled if needed, prioritizing lecture/tutorial time as everyone should be 
available then. Each meeting will be held according to the following standards:
\begin{itemize}
\item Each meeting will have a corresponding issue created, containing a previously made agenda, all decisions made, tasks assigned, deadlines and next steps/action items and anything else discussed
\item In case of member being unable to attend, they must Join a Microsoft Teams meeting to participate virtually if possible. If completely busy at that time, they must inform the rest of the team ~1 day prior if possible, as well as reviewing the agenda and providing any information relevant to the meeting, and then review any documents produced after
\item Roles will be assigned to team members to keep meetings structured and important items and tasks will not be missed
\end{itemize}

\section{Team Communication Plan}

There are many different channels of communication our team will employ. Communication channels include but are not limited to:
\begin{itemize}
	\item Primarily Microsoft Teams, split up into different channels to keep discussion organized
	\item Meetings and informal conversation between team members
	\item GitHub Issues
\end{itemize}
Any important information and decisions pertaining to the project should be recorded under the related GitHub issue or have a new one created if one does not exist already.

\section{Team Member Roles}

Roles are broken down into 3 sections: Meeting roles, Technical roles, and Research roles. Each team member should embrace all roles present, not leaving the responsibility of each of these roles entirely on one person, however their main focus should be on the roles assigned to them. Roles are subject to change in the future as well as being reassigned if the team sees fit.
\newline

\textbf{Meeting Roles:}
\begin{itemize}
\item Chair - Generally runs meetings, taking lead on agenda creation and keeping the team focused to cover all action items during meeting
\item Scribe - Records everything discussed, updating tasks, deadlines, and next steps/action items on the GitHub Repository
\item Progress Tracker - Follows up on details from previous meetings and deadlines to ensure team stays on track with older tasks
\item Quality/Scope Assessor - Assesses the scope of each decision made and ensures that team stays within the desired workload
\item Risk Assessor - Should question all assumptions made and attempt to poke holes in all decisions, should try and get team to think deeply about everything
\item Retrospective lead - Review what we have been doing and figure out what is and isn’t working for our team
\end{itemize}

\textbf{Technical Roles:}
\begin{itemize}
\item Systems Architect: takes lead on architecture design and oversees development is according to the blueprint
\item Back-end Engineer: Develops the back-end of the program
\item Front-end Engineer: Develops the front-end of the program
\item CI/CD Engineer: Focuses on creating workflows to support CI and CD
\end{itemize}

\textbf{Research Roles:}
\begin{itemize}
\item Research Design Lead: Formulates the research questions, hypotheses, and experimental methodology for evaluating the user experience.
\item Ethics and Participant Coordinator: Drafts consent forms and helps recruit study participants
\item Data Collection Analyst: Builds logging mechanisms into the software to capture quantitative user metrics (time-played, score, etc.)
\item Qualitative Research Lead: Conducts/designs post-test interviews, surveys and focus groups
\item Statistical Data Analyst: Analyzes quantitative data to drive further decisions
\item Visualization Specialist: Takes all data collected and converts it into an easy to understand format
\end{itemize}


\section{Workflow Plan}

\textbf{Git Development Workflow:} The project will be version controlled 
through the use of a GitHub repository, with reviewed and tested source code 
and documentation stored in the `main' branch. The project's files will be 
managed using the Feature Branch Git Workflow. The complete planned workflow is 
described below. When working on source code or documentation, team members 
will:

\begin{itemize}
	\item Start from the main branch, and ensure their local copy is up to date.
	\item Create a new Git feature branch to push their work to.
	\item Complete their work, and regularly stage and commit changes to the 
	feature branch.
	\item Push their changes to the feature branch.
	\item Create a pull request to merge their changes with the main branch.
	\item Request fellow team members to review their pull request.
	\item Merge the pull request with the main branch, once other team members 
	have completed their review and approved the request.
\end{itemize}

\noindent \textbf{Branches:} Branches should have descriptive names, related to 
the work being done. They should be written with lower case text, and multiple 
words should be hyphenated.
\newline \newline
\noindent \textbf{Pull Requests:} Pull Requests (PRs) should follow the 
guidelines below:

\begin{itemize}
	\item PRs should have descriptive names, detailing the feature or document 
	section being updated, and what was changed.
	\item PRs should have detailed descriptions, explaining the work that was 
	done in their changes, and what must be reviewed.
	\item PRs should be linked to GitHub issues that are related to their 
	changes, such as the feature(s) or documentation section(s) that were 
	completed through the changes.
	\item PRs should state the reviewer(s) that are planned to approve them. 
	Reviewers can be previously planned when beginning work on the issue, or 
	selected when the changes are planned to be merged (after discussion with 
	the team). Aside from these specific reviewers, all team members are 
	permitted to provide feedback on PRs, if necessary.
	\item Team members that created a PR and worked on its content, as well as 
	planned reviewers, should be assigned to the PR on GitHub.
\end{itemize}

\noindent \textbf{Issues:} Issues should follow the guidelines below:

\begin{itemize}
	\item Issues should have descriptive, title case names, detailing the work 
	product involved (e.g.~`Feature A'). If needed, issue names can also 
	contain descriptions of what will be done to complete them (e.g.~`Build 
	System B', or `Fix System C Bug'). For documentation, issues pertaining to 
	a specific document section should be named using the format `Document: 
	Section'.
	\item Issues should have detailed descriptions, clearly stating the work 
	required to complete each issue.
	\item Issues should have their task/issue owner assigned to them on GitHub.
	\item Sub issues will not be used, to minimize unneeded overhead.
	\item Documentation focused issues should use the Documentation issue 
	template.
	\item Source code feature focused issues should use the Feature issue 
	template.
	\item Bug fix focused issues should use the Bug issue template.
	\item Issues may use other labels and priorities as needed.
	\item Issues should be linked to relevant Projects and Milestones.
\end{itemize}

\noindent \textbf{Checklists:} Provided deliverable checklists should be 
reviewed by all team members before any required course deliverables are 
submitted.
\newline \newline
\noindent \textbf{Milestones:} Milestones should be created for each major 
course deliverable. They should have a descriptive name, and should be linked 
to the 
issues required to complete the deliverable.
\newline \newline
\noindent \textbf{CI/CD:} CI/CD pipelines will be used in the project's 
repository, through GitHub Actions. At a baseline level, CI/CD will be used for 
automated tests of source code, merged into the main project repository branch. 
Use of CI/CD pipelines may be expanded to pull request testing, code linting, 
or other purposes if needed (and if time is available for their construction).



\section{Project Decomposition and Scheduling}

GitHub projects will be used to keep track of upcoming course/project 
deliverables (as milestones), and the status of work tasks (as issues) linked 
to those deliverables. The GitHub project used for our team's repository is 
linked here: 
\href{https://github.com/orgs/Aeolus-Software/projects/2}{{RealmForge 
Project Tracker}}. This project contains a primary issue backlog 
table view, with fields displaying each issue's task owner (assignee), its 
status (`Todo', `In progress', `In review', or `Done'), any linked pull 
requests, the milestone the issue contributes towards, its type (Task, Feature, 
or Bug), and its target date (matching the deadline of the milestone it is 
linked to).
\newline \newline
When work begins for an upcoming 
milestone, all team members will collaborate in recurring team meetings to 
divide the milestone into 
suitable work tasks, which will be created as issues and linked to the GitHub 
project. This will involve splitting documentation by different sections, and 
subjects. For source code, development will be split across natural 
demarcations (e.g.~split into `Feature A' and `Feature B', or `Test System C' 
and `Test System D'). Reviewers will be selected for each issue, and each team 
member will complete their assigned tasks, collaborating with other members as 
needed. The division of work will be reassessed for a given milestone if any 
team members suggest that they unable to complete their work without assistance.
\newline \newline
The initial project schedule below corresponds to GitHub milestones created and 
linked to the GitHub project mentioned above. This schedule may be revised if 
any deliverable logistics change, or 
if the team chooses to alter it following more detailed planning (i.e.~pushing 
due dates earlier to prepare earlier, etc.).
\newline \newline
\noindent Initial project milestone schedule:

\begin{itemize}
	\item M1 Problem Statement, Development Plan, PoC Plan, and Team Charter 
	(Due by September 28, 2026)
	\item M2 SRS and Hazard Analysis Revision 0 (Due by October 19, 2026)
	\item M3 VnV Plan Revision 0 (Due by November 2, 2026)
	\item M4 Design Document Revision -1 and CD-Fall (FODA) (Due by November 
	16, 2026)
	\item M5 PoC Demo (Due by November 23, 2026)
	\item M6 Design Documentation Revision 0 (Due by January 18, 2027)
	\item M7 Revision 0 Demo (Due by February 1, 2027)
	\item M8 VnV Report Revision 0 and CD-Winter Revision 0 (Usability Report) 
	(Due by March 8, 2027)
	\item M9 Final Demonstration Revision 1, EXPO Demonstration (Due by March 
	15, 2027)
	\item M10 Final Documentation (Revision 1) (Due by March 31, 2027)
\end{itemize}

\section{Proof of Concept Demonstration Plan}

The main risks the team is worried about are an initial \textbf{Learning Curve}, the possibility of \textbf{Scope Creep}, and the resulting implementation being \textbf{Hard to Modify/Vary}.
\\

\noindent With whatever expected technology the team will use, there will certainly be a learning curve for the game engine/framework chosen. This could prove to be a barrier when attempting to create more complex portions of the project, and can make it hard to write robust code stright away since this is new software. 
\\ 

\noindent Scope creep is something the team sees as a risk as well. With today's technologies such as AI and LLMs it is much easier to implement new features without thinking it all the way through. If scope creep becomes an issue it may become more and more difficult to implement new features, since building features without solidifying the core components of the project leads to them being more complex and taking more time to implement. 
\\

\noindent The aforementioned risks culminate into a weak code base, which could be less modifiable and extendable than the project demands it to be. This would make implementing the myriad of feature variations a very difficult task to the point of not being able to meet deadlines.
\\

\noindent To prove these risks can be overcome, the proof of concept demonstration will include strong core components. To avoid scope creep it will stick to these core components so that they are very robust and extendable by the time they are finished. This means:
\begin{itemize}
        \item Basic unit testing with good code coverage
        \item Traversal to different areas/encounters through variable branched paths
        \item The main encounter/combat loop, with features like:
        \begin{itemize}
                 \item Turns/Turn order
                 \item Player/Enemy properties affecting gameplay
                 \item Items/Skills affecting Player/Enemy properties
                 \item Encounter completion conditions and reward options (different between runs)
        \end{itemize}
\end{itemize}


\noindent Important gameplay features will have elements of \textbf{modifiability} and \textbf{scalability}. This will ensure that the features are robust and that the current \textbf{feature variability} demonstrated sets a good precedent for implementing further variability, which won't hinder the project more than the team can handle.
\\

\noindent In the event of a failed proof of concept, the team has various methods of regrouping. First, the project's timeline could be reworked to have more rigid plans and deadlines for features. This would aid the direction of the project and ensure that it stays on course, avoiding scope creep through adding unnecessary features. Another method to refocus could be pivoting some key structure of the game to make it easier to fully flesh out. For example, changing the game engine/framework to something more familiar would allow the team to worry more about the quality of the work rather than learning while developing features. 


%\begin{itemize}
%    \item What will you demonstrate?
%	    \begin{itemize}
%	        \item The key components to our game
%	        \item The player is given the ability to traverse to different forking areas to loot
%	        \item The main combat loop (how the interacts and defeats the encounter)
%	        \item The player is given the ability to recieve loot (optional: given a choice to pick which to loot to take)  
%	    \end{itemize}
%    \begin{itemize}
%        \item \textbf{Inexperience with the technologies} and framework being used which leads to us unable to finish key %components of our game
%        \item All of the group members need to learn the coding framework and other technologies and be familiar with them in the %allotted timeframe
%        \item \textbf{Unable to create viable unit testing mechanisms}
%        \item Difficulting in utilizing the fullest code coverage tools for our framework
%        \item If our software testing is not proficient enough then certain bugs and logic errors might fall through delaying our %progress
%    \end{itemize}
%
%\end{itemize}
\section{Expected Technology}

%\wss{What programming language or languages do you expect to use?  What external
%libraries?  What frameworks?  What technologies.  Are there major components of
%the implementation that you expect you will implement, despite the existence of
%libraries that provide the required functionality.  For projects with machine
%learning, will you use pre-trained models, or be training your own model?  }

%\wss{The implementation decisions can, and likely will, change over the course
%of the project.  The initial documentation should be written in an abstract way;
%it should be agnostic of the implementation choices, unless the implementation
%choices are project constraints.  However, recording our initial thoughts on
%implementation helps understand the challenge level and feasibility of a
%project.  It may also help with early identification of areas where project
%members will need to augment their training.}

%Topics to discuss include the following:

% \begin{itemize}
% \item Specific programming language
% \item Specific libraries
% \item Pre-trained models
% \item Specific linter tool (if appropriate)
% \item Specific unit testing framework
% \item Investigation of code coverage measuring tools
% \item Specific plans for Continuous Integration (CI), or an explanation that CI
%   is not being done
% \item Specific performance measuring tools (like Valgrind), if
%   appropriate
% \item Tools you will likely be using?
% \end{itemize}

%\wss{git, GitHub and GitHub projects should be part of your technology.}

There are a number of game engines and frameworks that vary between ease of use, flexibility, cross-platform availability and other capabilities.
The team has narrowed it down to mainly Godot, GameMaker, Unity and LibGdx based on their respective pros and cons, although more programming technologies can be postulated if those prove not-suitable to the task. 
If the project uses Godot or GameMaker then the team follows the engine's respective language GDScript or GameMakerLanguage(GML), similarly if the project uses Unity or LibGDX then the strongly typed languages C\# or Java will be used.
\newline
\newline
The project will then have access to the above engine/framework's libraries and built-in APIs for use. The linter tools, code-coverage, and unit-test modules will once again vary depending on the engine. For example, if the libGDX framework is
is used, then the project can use the combination of JUnit (unit testing), JaCoCo (code coverage) and Spotbugs+Checkstyle (linting) for continuous integration/development (CI/CD). As for Godot, GdUnit4, GdUnit4 coverage, and GdLint are 
available and so on so forth with the others.
\newline
\newline
The project will utilize GitHub Actions for CI/CD after installation of the code testing modules to ensure proper bugs and errors are caught immediately.
A team-member will have the luxury to work on a sub-branch and perform pull requests and pushes without having to worry about breaking another teammates codebase.
This will trigger GitHub actions to set-up the framework environment, build the project, and perform all specified unit tests and code coverage. If the push or pull request fails one of the code coverage or unit tests, then the push is rejected.
Git and GitHub projects will be used for proper version-control of the code and ease of access via the cloud storage.
The team will then implement the roguelike functionality instead of using previous implements (if any) as well as the main gameplay loop in order to have full control and familiarity over manipulating game mechanics, game variables and other processes.

\section{Use of AI and LLMs}

Throughout this project, the team will use LLMs (Large Language Models) as 
needed to find/confirm common code syntax, grammatical/spelling rules, and 
other minor clarifications. LLMs may also be used to conduct brief `sanity 
checks' of documentation and source code. Finally, if determined to be 
beneficial and appropriate by a team member, LLMs may be used to generate 
source code, primarily for small sections and methods (such as implementing a 
set of functions to request an object's basic, primary attributes). The use of 
LLMs to generate large sections of source code will be avoided, and must be 
discussed with the team if a team member intends to use an LLM for this 
purpose. LLMs will not be used to generate documentation text, beyond simple 
automations or grammatical edits that do not significantly change a document's 
presented ideas (i.e.~creating an itemized list of text that was already 
written by a team member, vs.~generating an entire document section).
\newline \newline
Any work products (documentation, source code) generated by LLMs must be 
reviewed and verified for correctness by the team member using the LLM, and the 
use of the LLM must be transparently disclosed in each work product where it 
was 
used. Team members will not use any work products generated by LLMs that they 
cannot explain themselves or fully understand.
\newline \newline
A variety of common free LLMs may be used throughout the project, 
such as OpenAI's ChatGPT, Anthropic's Claude, Google Gemini, or Microsoft 
Copilot. Other premium versions of these tools may be used if team members 
decide to purchase personal access.

\section{Coding Standard}

The coding standard used throughout the project will depend on the game engine 
or framework used to develop the game:

\begin{itemize}
	\item If Godot is selected, the Godot documentation's GDScript style guide 
	\cite{GodotGDStyleGuide} will be used as a standard. If Godot's option for 
	writing code in C\# is used, the Godot documentation's C\# style guide 
	\cite{GodotCSharpStyleGuide} will 
	be used, to follow an appropriate style for C\# development in a video game 
	context.
\end{itemize} 

\newpage{}

\section{Appendix --- Generative AI Use Disclosure}

\subsection{AI Tools Used}

\begin{itemize}
	\item OpenAI GPT-5.6 Sol
	\item Google Search AI Overview (Google Gemini)
\end{itemize}

\subsection{How and Where Used}

AI tools were used for assistance with LaTeX syntax and commands, during the 
development of the Workflow Plan, Project Decomposition and Scheduling, Use of 
AI and LLMs, and Coding Standard sections. AI tools were not used to generate 
any text or ideas presented in these sections.

\subsection{Verification of AI-Generated Content}

Any LaTeX syntax provided by AI tools was verified by testing compilation in 
LaTeX documents, and testing the resulting output. Text was not generated by AI 
tools, and therefore no AI-generated text was required to be verified.

\newpage{}

\section*{Appendix --- Reflection}

\wss{Not required for CAS 741}

\input{../Reflection.text}

\begin{enumerate}
    \item Why is it important to create a development plan prior to starting the
    project?
    \item In your opinion, what are the advantages and disadvantages of using
    CI/CD?
    \item What disagreements did your group have in this deliverable, if any,
    and how did you resolve them?
\end{enumerate}

\newpage{}

\section*{Appendix --- Team Charter}

\subsection*{External Goals}

The team’s goal for this project is to aim for very high marks in this class and to produce a project we can all be proud to present at the end of the year, possibly continuing development after if desired. We are not aiming to win a prize at the Capstone EXPO.

\subsection*{Attendance}

\subsubsection*{Expectations}

Members should attend all meetings if possible, but if they need to miss a meeting, arrive later, or leave early, it is acceptable as long as they notify the team prior and provide any necessary details pertaining to the meeting.

\subsubsection*{Acceptable Excuse}

Having a class during the time of the meeting, having a meeting arranged earlier with another team, needing to complete work with a sooner deadline than our project, and personal affairs that meet a certain level of seriousness are all acceptable excuses to miss meetings as long as all their work is still completed. Excuses which are not acceptable include working on less urgent classwork, because they simply don’t think it is important or forgetting about the meeting.

\subsubsection*{In Case of Emergency}

The team member should notify the rest of the team as soon as they can that something is going on that is hindering them from doing work or attending a meeting. If possible, they should also send all relevant information they have to the team relating to the task they are no longer able to complete. We understand that emergencies happen that can stop members from completing tasks, but as long as they take reasonable steps to allow other team members to carry on without them and split up their work easily, the team should be able to manage without that team member for the period of their emergency.

\subsection*{Accountability and Teamwork}

\subsubsection*{Quality}

Team members should be aware of upcoming deliverables and their assigned responsibilities, and come to meetings ready for open discussions. Team members are encouraged to come prepared with questions or points of confusion so they can be discussed and resolved together.    

Deliverables should meet the requirements and be high quality, with no grammatical errors and accurate information that has been shared with the group.  

\subsubsection*{Attitude}

Team members should bring a positive attitude to all things work-related and be generally accepting and open-minded towards other team members and their ideas. We should not need a code of conduct and just use common sense to be respectful people to each other. Team members should only be harsh on each other if they are not meeting the team’s expectations in any way.

\subsubsection*{Stay on Track}

Weekly recurring team meetings will be scheduled to ensure the team stays on track. During these meetings, team members will share their updates and set goals for the upcoming week. By doing so, all members will be aware of the project's progression, the collective goals of the team, and the next steps being worked towards. 
 
Team members performing well will be recognized through the weekly meetings, as each member's contributions will be reviewed to ensure the group's quality of work is maintained. For team members who are underperforming, the group will hold open conversations to assess how the team can operate differently. This may be identified through repeated missed meetings and incomplete assigned responsibilities. If the behaviour continues, the team will discuss the issue with the TA.

\subsubsection*{Team Building}

The team will not have any specific team-building activities, but will be expected to bond simply through working together on a project we enjoy.

\subsubsection*{Decision Making}

Decisions the team makes should mostly be compromises of all team members' 
opinions and ideas until we can come to a consensus that we all agree with the 
idea. If it is not possible to compromise or come to a consensus, we will have 
to vote on all the ideas we have presented so far.

\newpage{}

\bibliographystyle{IEEEtran}
\bibliography{../../refs/References}

\end{document}
